\section{Introduction}
\label{sec:intro}

Vietnamese visual question answering (VQA) has a small number of strong open
models. \textbf{Vintern-1B}~\cite{vintern1b} is representative: a
1-billion-parameter vision--language model (VLM) that pairs an
InternViT-300M~\cite{internvl} encoder with a Qwen2-0.5B~\cite{qwen2} decoder
through a two-layer MLP projector, and reaches state-of-the-art accuracy on
Vietnamese VQA benchmarks. Its strength comes at a cost, however. Adapting
Vintern-1B to a new benchmark, as its authors do, means fully fine-tuning the
vision encoder and the projector and applying LoRA to the decoder, over roughly
three million image--question pairs. A recent alternative,
\textbf{ViMoE-VQA}~\cite{vimoevqa}, improves accuracy on the
AutoViVQA~\cite{autovivqa} benchmark by discarding this backbone entirely and
training a new mixture-of-experts VLM from scratch. Both routes are expensive,
and both rewrite most of the model.

This paper asks a narrower, more practical question:

\begin{quote}
\emph{Can we improve Vintern-1B on AutoViVQA by training only a small fraction of
its parameters --- freezing the entire backbone --- and if that is not enough,
where exactly is the bottleneck?}
\end{quote}

Concretely, we freeze both the InternViT-300M encoder and the Qwen2-0.5B decoder
and replace the original projector with a small trainable \emph{bridge} module
(0.78\,\% of all parameters). This is the cheapest possible adaptation: no
gradients flow into either backbone, training fits on a single 16\,GB GPU, and
one run takes a few hours rather than days. We then treat the remaining
performance gap as a diagnostic target. There are two places one could spend a
small additional parameter budget --- the \textbf{vision side} (a larger bridge,
more image tiles, adaptive per-question visual compute) or the \textbf{language
side} (a richer training signal, representation alignment to the decoder, or a
light adapter on the decoder itself) --- and we work through six research
questions, one lever at a time, to find which of them actually moves the needle.

\paragraph{Findings.}
\textbf{The recipe works, cheaply.} A frozen-backbone multi-token bridge at one
tile already surpasses the fully fine-tuned Vintern-1B on BLEU ($+9.4$), METEOR
($+5.0$) and CIDEr ($+23.7$), and beats ViMoE-VQA on corpus CIDEr-D, BLEU-4 and
ROUGE-L --- at about 1\,\% of the trainable parameters and no backbone
fine-tuning. A rank-16 LoRA on the decoder's attention projections (a further
$0.23\%$) lifts token-F1 from $49.6$ to $54.7$ and puts the recipe above
Vintern-1B on \emph{every} generation metric. Held-out test numbers match
validation within $0.5$ F1.

\textbf{The bottleneck is the decoder's attention, and nothing on the vision
side.} Of the six interventions, five are on the vision or training side --- a
$10\times$-larger bridge, more image tiles, a learned per-question routing policy,
multi-reference training, and projector-level representation alignment --- and all
five leave token-F1 unchanged or worse (alignment is an absolute null,
\dF{}~$-0.03$). The one that helps is a LoRA adapter on the frozen decoder; it
helps on \emph{every} bridge architecture, with a larger effect on weaker
bridges, after which all five collapse into a $0.9$-point F1 band. Moving the same
budget to the decoder's feed-forward layers instead diverges training. The useful
headroom is specifically in the frozen decoder's attention. This is also a direct
counterpoint to ViMoE-VQA's ``reasoning-aware'' account: on the same benchmark,
question type carries no usable signal about how much visual computation a
question needs.

\paragraph{Contributions.}
\begin{enumerate}
\item \textbf{A parameter-efficient adaptation recipe for Vietnamese VLMs}
  (Sec.~\ref{sec:method}, Sec.~\ref{sec:results}): a frozen InternViT + frozen
  Qwen2 + a lightweight pooled bridge + a rank-16 attention LoRA on the decoder
  --- ${\sim}1\%$ trainable parameters, one 16\,GB GPU, no backbone fine-tuning
  --- that matches or beats a fully fine-tuned Vintern-1B on generation metrics
  and holds up on a held-out test split.
\item \textbf{A systematic bottleneck diagnosis} (Sec.~\ref{sec:ablation}): a
  six-question ablation ladder over the vision and language sides of a
  frozen-backbone VLM. Four independent vision/training-side levers produce no
  lift; only decoder-attention capacity does, on every bridge. We localise the
  ceiling to the frozen decoder's attention projections --- feed-forward LoRA at
  the same budget diverges.
\item \textbf{A reliable evaluation protocol} for AutoViVQA: a leak-free grouped
  70/15/15 split (no image shared across splits), 3-seed means with bootstrap
  confidence intervals, matched validation/test reporting, a human-judgment
  sanity check on the token-F1 metric, and a FLOPs/latency/throughput
  characterisation of the image-tile lever that prior work on this benchmark
  explicitly deferred.
\end{enumerate}
